\documentclass[11pt]{article}
\usepackage[a4paper,margin=2cm]{geometry}
\usepackage{amsmath,amssymb}
\usepackage[hidelinks]{hyperref}
\allowdisplaybreaks
\title{Qpu.Processing.QuantumAlgorithm}
\author{Tsvetan Rouschev}
\date{@uuidna/qpu v1.1.1, CC-BY-NC-ND-4.0}
\begin{document}
\maketitle
\noindent\url{https://github.com/uuidna/qpu/blob/v1.1.1/lean/Qpu/Processing/QuantumAlgorithm.lean}

\paragraph{\texttt{QuantumAlgorithm.deutsch\_jozsa}}
one query decides: a constant function leaves all weight on 00, a balanced one leaves none
\[ \forall f < 16, (\mathrm{QuantumAlgorithm.constant}\ f = \mathrm{true} \to (\mathrm{QuantumAlgorithm.dj}\ f\ 0)\mathrm{.normSq} = 16) \land (\mathrm{QuantumAlgorithm.balanced}\ f = \mathrm{true} \to \mathrm{QuantumAlgorithm.dj}\ f\ 0 = 0) \]
\noindent axioms: none; proof: \url{https://github.com/uuidna/qpu/blob/v1.1.1/lean/Qpu/Processing/QuantumAlgorithm.lean\#L28}

\paragraph{\texttt{QuantumAlgorithm.grover\_search}}
one Grover iteration over four items leaves all weight on the marked item
\[ \forall m < 4, \forall r < 4, \mathrm{QuantumAlgorithm.grover}\ m\ r \ne 0 \leftrightarrow r = m \]
\noindent axioms: none; proof: \url{https://github.com/uuidna/qpu/blob/v1.1.1/lean/Qpu/Processing/QuantumAlgorithm.lean\#L39}

\paragraph{\texttt{QuantumAlgorithm.phase\_estimation\_exact}}
phase estimation: the register holding ω\^{}(jk) on |j⟩ for the eigenphase k/8 returns exactly |k⟩ under the inverse QFT
\[ \forall k < 8, \forall r < 8, (\mathrm{QuantumAlgorithm.sum8}\ \lambda \ j \mapsto \mathrm{QuantumAlgorithm.iqft}\ r\ j \cdot \mathrm{QuantumAlgorithm.wPow} (j \cdot k)) = \mathrm{if}\ r = k\ \mathrm{then}\ 8\ \mathrm{else}\ 0 \]
\noindent axioms: none; proof: \url{https://github.com/uuidna/qpu/blob/v1.1.1/lean/Qpu/Processing/QuantumAlgorithm.lean\#L53}

\paragraph{\texttt{QuantumAlgorithm.qft\_unitary}}
the QFT on three qubits is unitary: QFT · QFT† = 8 · I
\[ \forall j < 8, \forall k < 8, (\mathrm{QuantumAlgorithm.sum8}\ \lambda \ i \mapsto \mathrm{QuantumAlgorithm.qft}\ j\ i \cdot \mathrm{QuantumAlgorithm.iqft}\ i\ k) = \mathrm{if}\ j = k\ \mathrm{then}\ 8\ \mathrm{else}\ 0 \]
\noindent axioms: none; proof: \url{https://github.com/uuidna/qpu/blob/v1.1.1/lean/Qpu/Processing/QuantumAlgorithm.lean\#L51}
\end{document}
